\subsection{Coarse-to-Fine Cascade}\label{sec:methodology-cascade}

The cascade connects consecutive resolution levels through three
mechanisms, each replacing a coarser design choice made by Medverse or
Iris.

\paragraph{Query prior.} Medverse~\cite{huMedverseUniversalModel2025a}
feeds an upsampled, perturbed ground-truth mask into an extra U-Net branch
at every step; Iris~\cite{gaoShowSegmentUniversal2025} uses no prior at
all. At level $\ell > 1$, we instead initialize the target's
label token directly from level $\ell{-}1$'s own prediction:
%
\begin{equation}
  z^{\text{lbl}}_{\ell,0} = \mathrm{LblEmbed}\big(\mathrm{Warp}_{\ell-1 \to \ell}(\hat{L}^t_{\ell-1})\big),
\end{equation}
%
where $\mathrm{Warp}$ resamples level $\ell{-}1$'s prediction onto level
$\ell$'s crop grid, correcting for the two levels' different centers,
field of view, and any augmentation applied independently to each. This
makes the previous level's belief part of the same token stream that
attends bi-axially with the image (\autoref{sec:methodology-fusion}), rather
than a separate side channel. \autoref{fig:query-prior} shows this flow
across two levels.

\begin{figure}[t]
  \centering
  \includegraphics[width=\linewidth]{arch_query_prior.pdf}
  \caption{Query image prior across two cascade levels, with the full
  model input/output flow. At level 1 the target has no real prior yet
  (dashed support-mean placeholder); at level 2 the target's second lane
  is a real crop of level 1's own prediction $\hat{L}^t_1$ -- the same
  lane the context mask always occupies.}
  \label{fig:query-prior}
\end{figure}
\todo{Draft diagram, ported from the conference draft; finalize once the number of levels $M$ is settled.}

\paragraph{Region-restricted fine processing.} Medverse and Iris re-scan
the entire volume with a sliding window at every step. Once a coarse level
has localized the structure of interest, we instead restrict the sliding
window at the next, finer level to the coarse prediction's foreground
region, thresholded and padded by a fixed margin. Compute at fine levels
then scales with the size of the structure rather than the size of the
input volume, at the cost of depending on the coarse level's localization
being approximately correct.

\paragraph{Cross-level training via register tokens.} Medverse and Iris
connect levels only through the rendered prediction feeding the next
step, with no gradient path between them. We add a small set of learned
register tokens $\mathbf{r}_\ell$ that participate in level $\ell$'s
attention alongside the image and label tokens; their final state is
carried into level $\ell{+}1$ via a learned projection and a
level-independent type embedding:
%
\begin{equation}
  \mathbf{r}_{\ell+1,0} = W_r\, \mathbf{r}_{\ell,\text{final}} + \mathbf{e}_{\ell+1}.
\end{equation}
%
Unlike the query prior, this path is not detached between levels: level
$\ell{+}1$'s loss can backpropagate into level $\ell$'s own attention,
giving the model a direct incentive to carry forward whatever helps the
next level -- e.g.\ uncertainty or cues about neighboring structures --
beyond what the rendered mask alone can encode.
